% from ega4-8-fr.tex
\oldpage[IV-3]{26}

Rappelons (I, 4.1.10) que $\mathfrak{Spr}(X)$, pour tout préschéma $X$, est un
ensemble \emph{ordonné} par la relation «~$Z$ est un sous-préschéma du
sous-préschéma $Y$~» qui s'écrit $Z\leq Y$. Les applications
$u_{\lambda\mu}^{-1}$ et $u_\lambda^{-1}$ sont \emph{croissantes} pour les
relations d'ordre correspondantes dans $\mathfrak{Spr}(S_\lambda)$,
$\mathfrak{Spr}(S_\mu)$, $\mathfrak{Spr}(S)$. En outre, on définit une relation
d'ordre dans l'ensemble $\varinjlim\mathfrak{Spr}(S_\lambda)$ en écrivant que
$\eta\leq\eta'$ pour deux éléments de cet ensemble lorsqu'il existe un
$\lambda$ et deux éléments $Y_\lambda$, $Y'_\lambda$ de
$\mathfrak{Spr}(S_\lambda)$, dont $\eta$ et $\eta'$ sont les images canoniques,
et qui sont tels que $Y_\lambda\leq Y'_\lambda$~; on vérifie aisément que cela
ne dépend pas des représentants $Y_\lambda$, $Y'_\lambda$ considérés, et que
l'on a bien ainsi une \emph{relation d'ordre}. Cela étant, le fait que les
$u_\lambda^{-1}$ sont croissantes entraîne aussitôt que l'application
canonique (8.6.2.1) est \emph{croissante}~; il en est de même évidemment de
(8.6.2.2) et (8.6.2.3).

\begin{proposition}[8.6.3]
\label{IV.8.6.3-fr}
Supposons $S_0$ quasi-compact (resp. quasi-compact et quasi-séparé). Alors les
applications (8.6.2.1), (8.6.2.2), (8.6.2.3) sont injectives (resp.
bijectives).
\end{proposition}

\begin{proof}
Compte tenu des remarques de (8.6.1), les assertions relatives à (8.6.2.3)
découlent de (8.5.11) appliqué à $\mathcal F_\lambda=\mathcal O_{S_\lambda}$~;
de même, les assertions relatives à (8.6.2.2) sont des cas particuliers de
(8.3.5) et (8.3.11), compte tenu de ce que les $S_\lambda$ et $S$ sont
quasi-compacts. Reste à considérer l'application (8.6.2.1). Prouvons d'abord
qu'elle est surjective lorsque $S_0$ est quasi-compact et quasi-séparé. Soit
$Z$ un sous-préschéma de $S$, de présentation finie sur $S$~; comme $S$ est
quasi-compact, il en est de même de $Z$, donc il existe un ouvert
quasi-compact $U$ de $S$ tel que $Z$ soit un sous-préschéma fermé de $U$, de
présentation finie sur $U$. Il existe alors un indice $\lambda$ et un ouvert
quasi-compact $U_\lambda$ de $S_\lambda$ tels que
$U=u_\lambda^{-1}(U_\lambda)$ (8.2.11)~; comme $S_\lambda$ est quasi-séparé, il
en est de même de $U_\lambda$ (I, 1.2.7), et par suite on peut se borner au cas
où $U=S$~; mais dans ce cas, on est ramené au fait que (8.6.2.3) est
surjective.

Enfin, pour voir que (8.6.2.1) est injective lorsque $S_0$ est quasi-compact,
il suffira de prouver le résultat plus précis suivant~:
\end{proof}

\begin{corollary}[8.6.3.1]
\label{IV.8.6.3.1-fr}
Supposons $S_0$ quasi-compact et soient $Z'_\lambda$, $Z''_\lambda$ deux
sous-préschémas de $S_\lambda$, de présentation finie sur $S_\lambda$. Pour
que $Z'=u_\lambda^{-1}(Z'_\lambda)$ soit majoré par
$Z''=u_\lambda^{-1}(Z''_\lambda)$ (I, 4.1.10), il faut et il suffit qu'il
existe $\mu\geq\lambda$ tel que
$Z'_\mu=u_{\lambda\mu}^{-1}(Z'_\lambda)$ soit majoré par
$Z''_\mu=u_{\lambda\mu}^{-1}(Z''_\lambda)$.
\end{corollary}

\begin{proof}
Il est trivial que la condition est suffisante. Pour voir qu'elle est
nécessaire, notons d'abord que les ensembles sous-jacents $Z'_\lambda$ et
$Z''_\lambda$ sont localement constructibles dans $S_\lambda$ par hypothèse
(I, 1.8.4), donc l'hypothèse entraîne l'existence d'un $\nu\geq\lambda$ tel
que $Z'_\nu\subset Z''_\nu$ (8.3.5)~; remplaçant $\lambda$ par $\nu$, on peut
donc déjà supposer que l'on a $Z'_\lambda\subset Z''_\lambda$. En outre, par
hypothèse (I, 1.6.1), les sous-espaces $Z'_\lambda$ et $Z''_\lambda$ de
$S_\lambda$ sont quasi-compacts~; pour tout point $x\in Z'_\lambda$, il y a un
voisinage ouvert quasi-compact $V(x)$ dans $S_\lambda$ tels que
$V(x)\cap Z'_\lambda$ et $V(x)\cap Z''_\lambda$ soient fermés dans $V(x)$. En
recouvrant $S_\lambda$ par un nombre fini de voisinages $V(x_i)$ on voit donc
qu'il y a un ouvert quasi-compact $U_\lambda$ de $S_\lambda$ contenant
$Z'_\lambda$ et tel que $Z'_\lambda$ et $Z''_\lambda\cap U_\lambda$ soient
fermés dans $U_\lambda$. Si l'on désigne par $Y_\lambda$ le sous-préschéma
induit par $Z''_\lambda$ sur $U_\lambda\cap Z''_\lambda$, il est clair qu'avec
les notations habituelles, $Y_\mu$ (resp. $Y$) est induit par $Z''_\mu$ sur
$U_\mu\cap Z''_\mu$ pour $\mu\geq\lambda$ (resp. par $Z''$ sur
$U\cap Z''$)~; en outre $Z'$ est majoré par $Y$ (I, 4.4.1), et comme il suffit
de prouver que $Z'_\mu$ est majoré par $Y_\mu$ pour $\mu$
